\documentclass[../../main.tex]{subfiles}

\begin{document}

\section{一致可积性和紧性: 一般的 Vitali 收敛定理}\label{section:ch05-一般Vitali}

前一章的 Vitali 收敛定理告诉我们, 若 $m(E)<\infty$, $\{f_{n}\}$ 在 $E$ 上一致可积且在 $E$ 上几乎处处逐点收敛于 $f$, 则 $f$ 在 $E$ 上是可积的且在积分号下取极限是合理的, 即
\begin{align}
\lim_{n\to\infty}\left[\int_{E}f_{n}\right]=\int_{E}\lim_{n\to\infty}f_{n}=\int_{E}f.\label{align:ch05-有限Vitali}
\end{align}
该定理要求 $E$ 有有限测度. 事实上, 对每个自然数 $n$, 定义 $f_{n}=\chi_{(n,n+1)}$ 与 $f\equiv 0$. 则 $\{f_{n}\}$ 在 $\mathbf{R}$ 上一致可积且在 $\mathbf{R}$ 上逐点收敛于 $f$. 然而,
\begin{align*}
\lim_{n\to\infty}\left[\int_{E}f_{n}\right]=1\neq 0=\int_{E}\lim_{n\to\infty}f_{n}=\int_{E}f.
\end{align*}
无穷测度集上可积函数的以下性质暗示着对定义在无穷测度集上的函数序列, 为了保证积分号下取极限的合理性, 必须有与一致可积性对应的额外性质.

\begin{proposition}\label{proposition:ch05-尾部小}
令 $f$ 在 $E$ 上可积. 则对每个 $\varepsilon>0$, 存在有限测度的集合 $E_{0}$, 使得
\begin{align*}
\int_{E\sim E_{0}}|f|<\varepsilon.
\end{align*}
\end{proposition}

\begin{proof}
令 $\varepsilon>0$. 非负函数 $|f|$ 在 $E$ 上是可积的. 根据非负函数的积分的定义, 在 $E$ 上存在有界可测函数 $g$, 它在 $E$ 的有限测度子集 $E_{0}$ 外消失, 使得 $0\leqslant g\leqslant|f|$ 且 $\int_{E}|f|-\int_{E}g<\varepsilon$. 因此, 根据积分的线性与在区域上的可加性,
\begin{align*}
\int_{E\sim E_{0}}|f|=\int_{E\sim E_{0}}[|f|-g]\leqslant\int_{E}[|f|-g]<\varepsilon.
\end{align*}
\end{proof}

\begin{definition}[紧函数族]\label{definition:ch05-紧}
$E$ 上的可测函数族 $\mathcal{F}$ 称为在 $E$ 上\textbf{是紧的}, 若对每个 $\varepsilon>0$, 存在 $E$ 的有限测度的子集 $E_{0}$, 使得
\begin{align*}
\text{对所有 }f\in\mathcal{F},\quad\int_{E\sim E_{0}}|f|<\varepsilon.
\end{align*}
\end{definition}

我们从前一章的命题\cref{proposition:ch04-绝对连续积分}推出, 若 $\mathcal{F}$ 是 $E$ 上一致可积且紧的函数族, 则 $\mathcal{F}$ 中的每个函数在 $E$ 上是可积的.

\begin{theorem}[Vitali 收敛定理]\label{theorem:ch05-一般Vitali}
令 $\{f_{n}\}$ 为 $E$ 上的函数序列, 它在 $E$ 上是一致可积与紧的. 假定在 $E$ 上 a.e. 逐点地 $\{f_{n}\}\to f$. 则 $f$ 在 $E$ 上可积且
\begin{align*}
\lim_{n\to\infty}\int_{E}f_{n}=\int_{E}f.
\end{align*}
\end{theorem}

\begin{proof}
令 $\varepsilon>0$. 根据序列 $\{f_{n}\}$ 在 $E$ 上的紧性, 存在 $E$ 的可测子集 $E_{0}$, 其具有有限测度且对所有 $n$,
\begin{align*}
\int_{E\sim E_{0}}|f_{n}|<\varepsilon/4.
\end{align*}
我们从 Fatou 引理推出 $\int_{E\sim E_{0}}|f|\leqslant\varepsilon/4$. 因此 $f$ 在 $E\sim E_{0}$ 上可积. 此外, 根据积分的线性与单调性,
\begin{align}
\text{对所有 }n,\quad\left|\int_{E\sim E_{0}}[f_{n}-f]\right|\leqslant\int_{E\sim E_{0}}|f_{n}|+\int_{E\sim E_{0}}|f|<\varepsilon/2.\label{align:ch05-尾部估计}
\end{align}
但 $E_{0}$ 有有限测度而 $\{f_{n}\}$ 在 $E_{0}$ 上一致可积. 因此, 根据定义有限测度域上的函数的 Vitali 收敛定理, $f$ 在 $E_{0}$ 上可积, 且我们可以选取指标 $N$ 使得对所有 $n\geqslant N$,
\begin{align}
\left|\int_{E_{0}}[f_{n}-f]\right|<\varepsilon/2.\label{align:ch05-有限集估计}
\end{align}
因此 $f$ 在 $E$ 上可积, 根据\eqref{align:ch05-尾部估计}和\eqref{align:ch05-有限集估计}, 对所有 $n\geqslant N$,
\begin{align*}
\left|\int_{E}[f_{n}-f]\right|<\varepsilon.
\end{align*}
证毕.
\end{proof}

我们将以下推论的证明留作练习.

\begin{corollary}\label{corollary:ch05-非负收敛}
令 $\{h_{n}\}$ 为 $E$ 上的非负可积函数序列. 假定对 $E$ 中的几乎所有 $x$, $\{h_{n}(x)\}\to 0$. 则
\begin{align*}
\lim_{n\to\infty}\int_{E}h_{n}=0\quad\text{当且仅当}\{h_{n}\}\text{在 }E\text{ 上是一致可积和紧的}.
\end{align*}
\end{corollary}

\begin{exercise}\label{exercise:ch05-01}
令 $\{f_{k}\}_{k=1}^{n}$ 为有限函数族, 该族的每个函数在 $E$ 上可积. 证明在 $E$ 上 $\{f_{k}\}_{k=1}^{n}$ 是一致可积与紧的.
\end{exercise}

\begin{exercise}\label{exercise:ch05-02}
证明推论\cref{corollary:ch05-非负收敛}.
\end{exercise}

\begin{exercise}\label{exercise:ch05-03}
令函数序列 $\{f_{n}\}$ 和 $\{g_{n}\}$ 在 $E$ 上是一致可积与紧的. 证明对任何 $\alpha$ 和 $\beta$, $\{\alpha f_{n}+\beta g_{n}\}$ 在 $E$ 上也是一致可积与紧的.
\end{exercise}

\begin{exercise}\label{exercise:ch05-04}
令 $\{f_{n}\}$ 为 $E$ 上的可测函数序列. 证明 $\{f_{n}\}$ 在 $E$ 上是一致可积与紧的当且仅当对每个 $\varepsilon>0$, 存在 $E$ 的具有有限测度的可测子集 $E_{0}$ 和 $\delta>0$, 使得对 $E$ 的每个可测子集 $A$ 和指标 $n$,
\begin{align*}
\text{若 }m(A\cap E_{0})<\delta,\quad\text{则}\int_{A}|f_{n}|<\varepsilon.
\end{align*}
\end{exercise}

\begin{exercise}\label{exercise:ch05-05}
令 $\{f_{n}\}$ 为 $\mathbf{R}$ 上的可积函数序列. 证明 $\{f_{n}\}$ 在 $\mathbf{R}$ 上是一致可积与紧的当且仅当对每个 $\varepsilon>0$, 存在正数 $r$ 和 $\delta$ 使得对 $\mathbf{R}$ 的每个开子集 $\mathcal{O}$ 和指标 $n$,
\begin{align*}
\text{若 }m(\mathcal{O}\cap(-r,r))<\delta,\quad\text{则}\int_{\mathcal{O}}|f_{n}|<\varepsilon.
\end{align*}
\end{exercise}

\end{document}